% Author response template (LaTeX)
% Minimum reporting checklist for generative materials design
%
% Usage (standalone):
%   pdflatex author_response_template.tex
%
% Usage (in a manuscript):
%   Copy the section(s) and tables below into your SI/Appendix, or
%   % Author response template (LaTeX)
% Minimum reporting checklist for generative materials design
%
% Usage (standalone):
%   pdflatex author_response_template.tex
%
% Usage (in a manuscript):
%   Copy the section(s) and tables below into your SI/Appendix, or
%   % Author response template (LaTeX)
% Minimum reporting checklist for generative materials design
%
% Usage (standalone):
%   pdflatex author_response_template.tex
%
% Usage (in a manuscript):
%   Copy the section(s) and tables below into your SI/Appendix, or
%   % Author response template (LaTeX)
% Minimum reporting checklist for generative materials design
%
% Usage (standalone):
%   pdflatex author_response_template.tex
%
% Usage (in a manuscript):
%   Copy the section(s) and tables below into your SI/Appendix, or
%   \input{templates/author_response_template.tex} and delete this preamble.
%
\documentclass[11pt]{article}

\usepackage[margin=25mm]{geometry}
\usepackage{array}
\usepackage{booktabs}
\usepackage{longtable}
\usepackage[hidelinks]{hyperref}

\newcolumntype{L}[1]{>{\raggedright\arraybackslash}p{#1}}
\setlength{\tabcolsep}{3pt}

\begin{document}

\section*{Author response to the minimum reporting checklist}

\noindent\textbf{Paper title:} \emph{(fill in)}\\
\textbf{Version (date):} \emph{(fill in)}\\
\textbf{Manuscript ID (optional):} \emph{(fill in)}\\
\textbf{Corresponding author:} \emph{(fill in)}\\
\textbf{Code or data release (optional):} \emph{(URL, DOI, or repository tag/commit)}\\

\vspace{2mm}
\noindent\textbf{Instructions.} For each checklist item ID, provide a one to three sentence response and point to the exact location where the information is reported (main text section, SI section, figure, table, or repository path/commit). When space is limited, keep a one line summary in the main text and place full details in the SI.

\vspace{4mm}
\subsection*{Required items (must report in main text)}

\setlength{\LTpre}{0pt}
\setlength{\LTpost}{0pt}

\begin{longtable}{@{}L{0.06\textwidth} L{0.20\textwidth} L{0.50\textwidth} L{0.17\textwidth}@{}}
\toprule
\textbf{ID} & \textbf{Checklist item} & \textbf{Response (one to three sentences)} & \textbf{Where reported}\\
\midrule
\endfirsthead
\toprule
\textbf{ID} & \textbf{Checklist item} & \textbf{Response (one to three sentences)} & \textbf{Where reported}\\
\midrule
\endhead
\midrule
\multicolumn{4}{r}{\small Continued on next page}\\
\endfoot
\bottomrule
\endlastfoot

\texttt{D1} & Data & \emph{(fill in; one to three sentences)} & \emph{(e.g., Main text Sec.~X; SI Sec.~Sx)}\\[1mm]
\texttt{S1} & Splits and leakage & \emph{(fill in; one to three sentences)} & \emph{(e.g., Main text Sec.~X; SI Sec.~Sx)}\\[1mm]
\texttt{M1} & Model specification & \emph{(fill in; one to three sentences)} & \emph{(e.g., Main text Sec.~X; SI Sec.~Sx)}\\[1mm]
\texttt{T1} & Training and inference & \emph{(fill in; one to three sentences)} & \emph{(e.g., Main text Sec.~X; SI Sec.~Sx)}\\[1mm]
\texttt{G1} & Generation protocol & \emph{(fill in; one to three sentences)} & \emph{(e.g., Main text Sec.~X; SI Sec.~Sx)}\\[1mm]
\texttt{E1} & Evaluation & \emph{(fill in; one to three sentences)} & \emph{(e.g., Main text Sec.~X; SI Sec.~Sx)}\\[1mm]
\texttt{B1} & Baselines and ablations & \emph{(fill in; one to three sentences)} & \emph{(e.g., Main text Sec.~X; SI Sec.~Sx)}\\[1mm]
\texttt{F1} & Feasibility (thermodynamic vs synthetic) & \emph{(fill in; one to three sentences)} & \emph{(e.g., Main text Sec.~X; SI Sec.~Sx)}\\[1mm]
\texttt{R1} & Reproducibility & \emph{(fill in; one to three sentences)} & \emph{(e.g., Main text Sec.~X; SI Sec.~Sx)}\\[1mm]

\end{longtable}

\subsection*{Optional items (recommended)}

\begin{longtable}{@{}L{0.06\textwidth} L{0.20\textwidth} L{0.50\textwidth} L{0.17\textwidth}@{}}
\toprule
\textbf{ID} & \textbf{Checklist item} & \textbf{Response (one to three sentences)} & \textbf{Where reported}\\
\midrule
\endfirsthead
\toprule
\textbf{ID} & \textbf{Checklist item} & \textbf{Response (one to three sentences)} & \textbf{Where reported}\\
\midrule
\endhead
\midrule
\multicolumn{4}{r}{\small Continued on next page}\\
\endfoot
\bottomrule
\endlastfoot

\texttt{CL1} & Closed loop or active learning & \emph{(fill in; one to three sentences)} & \emph{(e.g., Main text Sec.~X; SI Sec.~Sx)}\\[1mm]
\texttt{C1} & Compute footprint & \emph{(fill in; one to three sentences)} & \emph{(e.g., Main text Sec.~X; SI Sec.~Sx)}\\[1mm]
\texttt{GOV1} & Governance and responsible use & \emph{(fill in; one to three sentences)} & \emph{(e.g., Main text Sec.~X; SI Sec.~Sx)}\\[1mm]
\texttt{K1} & Key number summary & \emph{(fill in; one to three sentences)} & \emph{(e.g., Main text Sec.~X; SI Sec.~Sx)}\\[1mm]

\end{longtable}

\end{document}
 and delete this preamble.
%
\documentclass[11pt]{article}

\usepackage[margin=25mm]{geometry}
\usepackage{array}
\usepackage{booktabs}
\usepackage{longtable}
\usepackage[hidelinks]{hyperref}

\newcolumntype{L}[1]{>{\raggedright\arraybackslash}p{#1}}
\setlength{\tabcolsep}{3pt}

\begin{document}

\section*{Author response to the minimum reporting checklist}

\noindent\textbf{Paper title:} \emph{(fill in)}\\
\textbf{Version (date):} \emph{(fill in)}\\
\textbf{Manuscript ID (optional):} \emph{(fill in)}\\
\textbf{Corresponding author:} \emph{(fill in)}\\
\textbf{Code or data release (optional):} \emph{(URL, DOI, or repository tag/commit)}\\

\vspace{2mm}
\noindent\textbf{Instructions.} For each checklist item ID, provide a one to three sentence response and point to the exact location where the information is reported (main text section, SI section, figure, table, or repository path/commit). When space is limited, keep a one line summary in the main text and place full details in the SI.

\vspace{4mm}
\subsection*{Required items (must report in main text)}

\setlength{\LTpre}{0pt}
\setlength{\LTpost}{0pt}

\begin{longtable}{@{}L{0.06\textwidth} L{0.20\textwidth} L{0.50\textwidth} L{0.17\textwidth}@{}}
\toprule
\textbf{ID} & \textbf{Checklist item} & \textbf{Response (one to three sentences)} & \textbf{Where reported}\\
\midrule
\endfirsthead
\toprule
\textbf{ID} & \textbf{Checklist item} & \textbf{Response (one to three sentences)} & \textbf{Where reported}\\
\midrule
\endhead
\midrule
\multicolumn{4}{r}{\small Continued on next page}\\
\endfoot
\bottomrule
\endlastfoot

\texttt{D1} & Data & \emph{(fill in; one to three sentences)} & \emph{(e.g., Main text Sec.~X; SI Sec.~Sx)}\\[1mm]
\texttt{S1} & Splits and leakage & \emph{(fill in; one to three sentences)} & \emph{(e.g., Main text Sec.~X; SI Sec.~Sx)}\\[1mm]
\texttt{M1} & Model specification & \emph{(fill in; one to three sentences)} & \emph{(e.g., Main text Sec.~X; SI Sec.~Sx)}\\[1mm]
\texttt{T1} & Training and inference & \emph{(fill in; one to three sentences)} & \emph{(e.g., Main text Sec.~X; SI Sec.~Sx)}\\[1mm]
\texttt{G1} & Generation protocol & \emph{(fill in; one to three sentences)} & \emph{(e.g., Main text Sec.~X; SI Sec.~Sx)}\\[1mm]
\texttt{E1} & Evaluation & \emph{(fill in; one to three sentences)} & \emph{(e.g., Main text Sec.~X; SI Sec.~Sx)}\\[1mm]
\texttt{B1} & Baselines and ablations & \emph{(fill in; one to three sentences)} & \emph{(e.g., Main text Sec.~X; SI Sec.~Sx)}\\[1mm]
\texttt{F1} & Feasibility (thermodynamic vs synthetic) & \emph{(fill in; one to three sentences)} & \emph{(e.g., Main text Sec.~X; SI Sec.~Sx)}\\[1mm]
\texttt{R1} & Reproducibility & \emph{(fill in; one to three sentences)} & \emph{(e.g., Main text Sec.~X; SI Sec.~Sx)}\\[1mm]

\end{longtable}

\subsection*{Optional items (recommended)}

\begin{longtable}{@{}L{0.06\textwidth} L{0.20\textwidth} L{0.50\textwidth} L{0.17\textwidth}@{}}
\toprule
\textbf{ID} & \textbf{Checklist item} & \textbf{Response (one to three sentences)} & \textbf{Where reported}\\
\midrule
\endfirsthead
\toprule
\textbf{ID} & \textbf{Checklist item} & \textbf{Response (one to three sentences)} & \textbf{Where reported}\\
\midrule
\endhead
\midrule
\multicolumn{4}{r}{\small Continued on next page}\\
\endfoot
\bottomrule
\endlastfoot

\texttt{CL1} & Closed loop or active learning & \emph{(fill in; one to three sentences)} & \emph{(e.g., Main text Sec.~X; SI Sec.~Sx)}\\[1mm]
\texttt{C1} & Compute footprint & \emph{(fill in; one to three sentences)} & \emph{(e.g., Main text Sec.~X; SI Sec.~Sx)}\\[1mm]
\texttt{GOV1} & Governance and responsible use & \emph{(fill in; one to three sentences)} & \emph{(e.g., Main text Sec.~X; SI Sec.~Sx)}\\[1mm]
\texttt{K1} & Key number summary & \emph{(fill in; one to three sentences)} & \emph{(e.g., Main text Sec.~X; SI Sec.~Sx)}\\[1mm]

\end{longtable}

\end{document}
 and delete this preamble.
%
\documentclass[11pt]{article}

\usepackage[margin=25mm]{geometry}
\usepackage{array}
\usepackage{booktabs}
\usepackage{longtable}
\usepackage[hidelinks]{hyperref}

\newcolumntype{L}[1]{>{\raggedright\arraybackslash}p{#1}}
\setlength{\tabcolsep}{3pt}

\begin{document}

\section*{Author response to the minimum reporting checklist}

\noindent\textbf{Paper title:} \emph{(fill in)}\\
\textbf{Version (date):} \emph{(fill in)}\\
\textbf{Manuscript ID (optional):} \emph{(fill in)}\\
\textbf{Corresponding author:} \emph{(fill in)}\\
\textbf{Code or data release (optional):} \emph{(URL, DOI, or repository tag/commit)}\\

\vspace{2mm}
\noindent\textbf{Instructions.} For each checklist item ID, provide a one to three sentence response and point to the exact location where the information is reported (main text section, SI section, figure, table, or repository path/commit). When space is limited, keep a one line summary in the main text and place full details in the SI.

\vspace{4mm}
\subsection*{Required items (must report in main text)}

\setlength{\LTpre}{0pt}
\setlength{\LTpost}{0pt}

\begin{longtable}{@{}L{0.06\textwidth} L{0.20\textwidth} L{0.50\textwidth} L{0.17\textwidth}@{}}
\toprule
\textbf{ID} & \textbf{Checklist item} & \textbf{Response (one to three sentences)} & \textbf{Where reported}\\
\midrule
\endfirsthead
\toprule
\textbf{ID} & \textbf{Checklist item} & \textbf{Response (one to three sentences)} & \textbf{Where reported}\\
\midrule
\endhead
\midrule
\multicolumn{4}{r}{\small Continued on next page}\\
\endfoot
\bottomrule
\endlastfoot

\texttt{D1} & Data & \emph{(fill in; one to three sentences)} & \emph{(e.g., Main text Sec.~X; SI Sec.~Sx)}\\[1mm]
\texttt{S1} & Splits and leakage & \emph{(fill in; one to three sentences)} & \emph{(e.g., Main text Sec.~X; SI Sec.~Sx)}\\[1mm]
\texttt{M1} & Model specification & \emph{(fill in; one to three sentences)} & \emph{(e.g., Main text Sec.~X; SI Sec.~Sx)}\\[1mm]
\texttt{T1} & Training and inference & \emph{(fill in; one to three sentences)} & \emph{(e.g., Main text Sec.~X; SI Sec.~Sx)}\\[1mm]
\texttt{G1} & Generation protocol & \emph{(fill in; one to three sentences)} & \emph{(e.g., Main text Sec.~X; SI Sec.~Sx)}\\[1mm]
\texttt{E1} & Evaluation & \emph{(fill in; one to three sentences)} & \emph{(e.g., Main text Sec.~X; SI Sec.~Sx)}\\[1mm]
\texttt{B1} & Baselines and ablations & \emph{(fill in; one to three sentences)} & \emph{(e.g., Main text Sec.~X; SI Sec.~Sx)}\\[1mm]
\texttt{F1} & Feasibility (thermodynamic vs synthetic) & \emph{(fill in; one to three sentences)} & \emph{(e.g., Main text Sec.~X; SI Sec.~Sx)}\\[1mm]
\texttt{R1} & Reproducibility & \emph{(fill in; one to three sentences)} & \emph{(e.g., Main text Sec.~X; SI Sec.~Sx)}\\[1mm]

\end{longtable}

\subsection*{Optional items (recommended)}

\begin{longtable}{@{}L{0.06\textwidth} L{0.20\textwidth} L{0.50\textwidth} L{0.17\textwidth}@{}}
\toprule
\textbf{ID} & \textbf{Checklist item} & \textbf{Response (one to three sentences)} & \textbf{Where reported}\\
\midrule
\endfirsthead
\toprule
\textbf{ID} & \textbf{Checklist item} & \textbf{Response (one to three sentences)} & \textbf{Where reported}\\
\midrule
\endhead
\midrule
\multicolumn{4}{r}{\small Continued on next page}\\
\endfoot
\bottomrule
\endlastfoot

\texttt{CL1} & Closed loop or active learning & \emph{(fill in; one to three sentences)} & \emph{(e.g., Main text Sec.~X; SI Sec.~Sx)}\\[1mm]
\texttt{C1} & Compute footprint & \emph{(fill in; one to three sentences)} & \emph{(e.g., Main text Sec.~X; SI Sec.~Sx)}\\[1mm]
\texttt{GOV1} & Governance and responsible use & \emph{(fill in; one to three sentences)} & \emph{(e.g., Main text Sec.~X; SI Sec.~Sx)}\\[1mm]
\texttt{K1} & Key number summary & \emph{(fill in; one to three sentences)} & \emph{(e.g., Main text Sec.~X; SI Sec.~Sx)}\\[1mm]

\end{longtable}

\end{document}
 and delete this preamble.
%
\documentclass[11pt]{article}

\usepackage[margin=25mm]{geometry}
\usepackage{array}
\usepackage{booktabs}
\usepackage{longtable}
\usepackage[hidelinks]{hyperref}

\newcolumntype{L}[1]{>{\raggedright\arraybackslash}p{#1}}
\setlength{\tabcolsep}{3pt}

\begin{document}

\section*{Author response to the minimum reporting checklist}

\noindent\textbf{Paper title:} \emph{(fill in)}\\
\textbf{Version (date):} \emph{(fill in)}\\
\textbf{Manuscript ID (optional):} \emph{(fill in)}\\
\textbf{Corresponding author:} \emph{(fill in)}\\
\textbf{Code or data release (optional):} \emph{(URL, DOI, or repository tag/commit)}\\

\vspace{2mm}
\noindent\textbf{Instructions.} For each checklist item ID, provide a one to three sentence response and point to the exact location where the information is reported (main text section, SI section, figure, table, or repository path/commit). When space is limited, keep a one line summary in the main text and place full details in the SI.

\vspace{4mm}
\subsection*{Required items (must report in main text)}

\setlength{\LTpre}{0pt}
\setlength{\LTpost}{0pt}

\begin{longtable}{@{}L{0.06\textwidth} L{0.20\textwidth} L{0.50\textwidth} L{0.17\textwidth}@{}}
\toprule
\textbf{ID} & \textbf{Checklist item} & \textbf{Response (one to three sentences)} & \textbf{Where reported}\\
\midrule
\endfirsthead
\toprule
\textbf{ID} & \textbf{Checklist item} & \textbf{Response (one to three sentences)} & \textbf{Where reported}\\
\midrule
\endhead
\midrule
\multicolumn{4}{r}{\small Continued on next page}\\
\endfoot
\bottomrule
\endlastfoot

\texttt{D1} & Data & \emph{(fill in; one to three sentences)} & \emph{(e.g., Main text Sec.~X; SI Sec.~Sx)}\\[1mm]
\texttt{S1} & Splits and leakage & \emph{(fill in; one to three sentences)} & \emph{(e.g., Main text Sec.~X; SI Sec.~Sx)}\\[1mm]
\texttt{M1} & Model specification & \emph{(fill in; one to three sentences)} & \emph{(e.g., Main text Sec.~X; SI Sec.~Sx)}\\[1mm]
\texttt{T1} & Training and inference & \emph{(fill in; one to three sentences)} & \emph{(e.g., Main text Sec.~X; SI Sec.~Sx)}\\[1mm]
\texttt{G1} & Generation protocol & \emph{(fill in; one to three sentences)} & \emph{(e.g., Main text Sec.~X; SI Sec.~Sx)}\\[1mm]
\texttt{E1} & Evaluation & \emph{(fill in; one to three sentences)} & \emph{(e.g., Main text Sec.~X; SI Sec.~Sx)}\\[1mm]
\texttt{B1} & Baselines and ablations & \emph{(fill in; one to three sentences)} & \emph{(e.g., Main text Sec.~X; SI Sec.~Sx)}\\[1mm]
\texttt{F1} & Feasibility (thermodynamic vs synthetic) & \emph{(fill in; one to three sentences)} & \emph{(e.g., Main text Sec.~X; SI Sec.~Sx)}\\[1mm]
\texttt{R1} & Reproducibility & \emph{(fill in; one to three sentences)} & \emph{(e.g., Main text Sec.~X; SI Sec.~Sx)}\\[1mm]

\end{longtable}

\subsection*{Optional items (recommended)}

\begin{longtable}{@{}L{0.06\textwidth} L{0.20\textwidth} L{0.50\textwidth} L{0.17\textwidth}@{}}
\toprule
\textbf{ID} & \textbf{Checklist item} & \textbf{Response (one to three sentences)} & \textbf{Where reported}\\
\midrule
\endfirsthead
\toprule
\textbf{ID} & \textbf{Checklist item} & \textbf{Response (one to three sentences)} & \textbf{Where reported}\\
\midrule
\endhead
\midrule
\multicolumn{4}{r}{\small Continued on next page}\\
\endfoot
\bottomrule
\endlastfoot

\texttt{CL1} & Closed loop or active learning & \emph{(fill in; one to three sentences)} & \emph{(e.g., Main text Sec.~X; SI Sec.~Sx)}\\[1mm]
\texttt{C1} & Compute footprint & \emph{(fill in; one to three sentences)} & \emph{(e.g., Main text Sec.~X; SI Sec.~Sx)}\\[1mm]
\texttt{GOV1} & Governance and responsible use & \emph{(fill in; one to three sentences)} & \emph{(e.g., Main text Sec.~X; SI Sec.~Sx)}\\[1mm]
\texttt{K1} & Key number summary & \emph{(fill in; one to three sentences)} & \emph{(e.g., Main text Sec.~X; SI Sec.~Sx)}\\[1mm]

\end{longtable}

\end{document}
